% !PRINTED 728
\weberpage{728}
%<<CONT
remaining ones, however, in reverse order and with altered sign.
In doing this $\frac{1}{2}\nu'$ changes of sign are necessary,
and we therefore obtain for the product (12) the determinant
defined by the following equation, whose absolute value we
denote by $T$:
\weq{14}{(-1)^{\frac{1}{2}\nu'}
\begin{vmatrix}
\lambda_0, & \lambda_1, & \ldots, & \lambda_{\nu'-2}, & \lambda_{\nu'-1}\\
\lambda_1, & \lambda_2, & \ldots, & \lambda_{\nu'-1}, & -\lambda_0\\
\lambda_2, & \lambda_3, & \ldots, & -\lambda_0, & -\lambda_1\\
\cdots & \cdots & \cdots & \cdots & \cdots\\
\lambda_{\nu'-1}, & -\lambda_0, & \ldots, & -\lambda_{\nu'-3}, & -\lambda_{\nu'-2}
\end{vmatrix} = \pm T,}
in which the negative signs stand to the right of the diagonal row
running from lower left to upper right.

The determinant $T$ is here so arranged that each row
arises from the preceding one by the substitution $(r, r^{5^\beta})$.
Each column therefore contains a system of conjugate logarithms.

By (12) we now have
\weq{15}{B = \frac{T}{D}.}
In order to gain clarity about the nature of this result,
it is necessary to consider somewhat more closely the number $D$, which was
defined by \S.\,195, (12), \S.\,197, (6); and this requires us to enter
 into the theory of the units of the field $H_m$.

\begin{center}
\S.\,200.\\
\gesperrt{Normal units in $H_m$.}
\end{center}

We now introduce a simplifying notation, by setting
\weq{1}{r^{5^\beta} = r_\beta}
and letting $\beta$ run through a complete system of residues with respect to the modulus $\nu$.
 From this it follows that
\weq{2}{r_{\beta+\nu'} = -r_\beta,}
and in the form $(r, r_\beta)$ not all the substitutions
of the group of the field $\Omega_m$ are then contained; rather, there
must still be added the substitutions $(r, r_\beta^{-1})$; but
$(r, r_\beta)$ does furnish us all the substitutions of the group of $H_m$.

We now denote by $\mathfrak{E}(r)$ a unit of the field $H_m$.
Among these are also the units of the field $H_\mu$ (as
%CONT>>
